\documentclass[10pt,a4paper]{amsart}
\usepackage{amsmath,amssymb,amsthm,graphicx}
\usepackage{verbatim}
\newenvironment{reference}{\comment}{\endcomment}
%\newenvironment{reference}{}{}
\newenvironment{slogan}{\comment}{\endcomment}
\newenvironment{history}{\comment}{\endcomment}

% For commutative diagrams we use Xy-pic
\usepackage[all]{xy}

% We use 2cell for 2-commutative diagrams.
\xyoption{2cell}
% Theorem environments.
%
\theoremstyle{plain}
\newtheorem{theorem}[subsection]{Theorem}
\newtheorem{proposition}[subsection]{Proposition}
\newtheorem{lemma}[subsection]{Lemma}

\theoremstyle{definition}
\newtheorem{definition}[subsection]{Definition}
\newtheorem{example}[subsection]{Example}
\newtheorem{exercise}[subsection]{Exercise}
\newtheorem{situation}[subsection]{Situation}

\theoremstyle{remark}
\newtheorem{remark}[subsection]{Remark}
\newtheorem{remarks}[subsection]{Remarks}

\numberwithin{equation}{subsection}

% Macros
%
\def\lim{\mathop{\mathrm{lim}}\nolimits}
\def\colim{\mathop{\mathrm{colim}}\nolimits}
\def\Spec{\mathop{\mathrm{Spec}}}
\def\Hom{\mathop{\mathrm{Hom}}\nolimits}
\def\Ext{\mathop{\mathrm{Ext}}\nolimits}
\def\SheafHom{\mathop{\mathcal{H}\!\mathit{om}}\nolimits}
\def\SheafExt{\mathop{\mathcal{E}\!\mathit{xt}}\nolimits}
\def\Sch{\mathit{Sch}}
\def\Mor{\mathop{\mathrm{Mor}}\nolimits}
\def\Ob{\mathop{\mathrm{Ob}}\nolimits}
\def\Sh{\mathop{\mathit{Sh}}\nolimits}
\def\NL{\mathop{N\!L}\nolimits}
\def\CH{\mathop{\mathrm{CH}}\nolimits}
\def\proetale{{pro\text{-}\acute{e}tale}}
\def\etale{{\acute{e}tale}}
\def\QCoh{\mathit{QCoh}}
\def\Ker{\mathop{\mathrm{Ker}}}
\def\Im{\mathop{\mathrm{Im}}}
\def\Coker{\mathop{\mathrm{Coker}}}
\def\Coim{\mathop{\mathrm{Coim}}}

% Boxtimes
%
\DeclareMathSymbol{\boxtimes}{\mathbin}{AMSa}{"02}

%
% Macros for moduli stacks/spaces
%
\def\QCohstack{\mathcal{QC}\!\mathit{oh}}
\def\Cohstack{\mathcal{C}\!\mathit{oh}}
\def\Spacesstack{\mathcal{S}\!\mathit{paces}}
\def\Quotfunctor{\mathrm{Quot}}
\def\Hilbfunctor{\mathrm{Hilb}}
\def\Curvesstack{\mathcal{C}\!\mathit{urves}}
\def\Polarizedstack{\mathcal{P}\!\mathit{olarized}}
\def\Complexesstack{\mathcal{C}\!\mathit{omplexes}}
% \Pic is the operator that assigns to X its picard group, usage \Pic(X)
% \Picardstack_{X/B} denotes the Picard stack of X over B
% \Picardfunctor_{X/B} denotes the Picard functor of X over B
\def\Pic{\mathop{\mathrm{Pic}}\nolimits}
\def\Picardstack{\mathcal{P}\!\mathit{ic}}
\def\Picardfunctor{\mathrm{Pic}}
\def\Deformationcategory{\mathcal{D}\!\mathit{ef}}
\usepackage{cleveref}
\crefname{originalreference}{original reference}{original references}
\Crefname{originalreference}{Original reference}{Original references}
\makeatletter
\newlabel{lemma-divide-difference-points}{{0FH7}{source}{}{sourcepointer-1}{}}
\newlabel{lemma-divide-difference-points@cref}{{[originalreference][0][]0FH7}{[1][0][]source}{}{}{}}
\makeatother
\crefname{originalreference}{original reference}{original references}
\Crefname{originalreference}{Original reference}{Original references}
\makeatletter
\newlabel{lemma-generated-by-separable}{{0FH5}{source}{}{sourcepointer-2}{}}
\newlabel{lemma-generated-by-separable@cref}{{[originalreference][0][]0FH5}{[1][0][]source}{}{}{}}
\makeatother
\crefname{originalreference}{original reference}{original references}
\Crefname{originalreference}{Original reference}{Original references}
\makeatletter
\newlabel{lemma-relations-classes-points}{{0FHU}{source}{}{sourcepointer-3}{}}
\newlabel{lemma-relations-classes-points@cref}{{[originalreference][0][]0FHU}{[1][0][]source}{}{}{}}
\makeatother
\title{The cycle class of a cycle of points depends only on its degree}
\author{The Stacks Project Authors}
\date{stacks-project:weil.tex:local-snapshot-20261012; GNU FDL 1.2 or later}
\begin{document}
\maketitle
\noindent\textit{Editorial note.} All delivered mathematical text is the author's own, copied from the pinned Stacks Project chapter file weil.tex; no statement or proof step was rewritten, reordered or removed.\par\medskip
\noindent\textit{Editorial note.} The counted proof is the author's complete argument: the reduction to points with separable residue fields, the base change along a finite Galois extension producing rational points, and the divisibility argument through the relations lemma.\par\medskip
\noindent\textit{Editorial note.} The label lemma-classes-points names the counted statement, and the proof environment immediately following it in the source is delivered verbatim as its proof.\par\medskip
\noindent\textit{Editorial note.} The counted statement is carried by the author's own lemma environment in the chapter source.\par\medskip
\noindent\textit{Editorial note.} Pre-open gate: python operations/agent-authoring/tools/precard\_gate.py --source weil.tex --label lemma-classes-points --start 3696 --end 3734 --local-id 28140 returned RESULT=OK (L28140 free, occupancy FREE, public interval clear). Pre-assembly tag self-check: canonical tag resolved and not present in reference/public-items.jsonl.\par\medskip
\noindent\textit{Editorial note.} The statement and proof are copied from the pinned chapter file; the cited inputs are kept as external pointers with their canonical Stacks tags.\par\medskip
\noindent\textit{Allowed standard prerequisites (editorial).}\begin{itemize}
\item Generation of the group of zero cycles by classes of points with separable residue field. Same chapter, the lemma that the group of zero cycles of a smooth projective scheme is generated by classes of closed points with separable residue field. Reduces the statement to a comparison of two closed points with separable residue fields.
\item Divisibility of the difference of two rational points after a separable extension. Same chapter, the lemma producing a finite separable extension over which the difference of two rational points becomes divisible by an integer larger than one. Produces the extension used to reduce the comparison of rational points to the case treated by the relations lemma.
\item Compatibility of the cycle class map with pushforwards along a finite morphism. Axiom (C) on the compatibility of the cycle class map with proper pushforwards, and the lemma relating the classes of points of the same degree. Transfers the equality of cycle classes from the base change back to the original scheme.
\end{itemize}\medskip
\ifdefined\hypertarget\hypertarget{sourcepointer-1}{}\fi
\noindent\textit{Source reference.} Original reference 0FH7: sources/stacks-weil/pinned-original-source.tex, source lines 2325--2325. Prerequisite pointer invoked at proof line(s) 3724; it names a standard fact of the same chapter that the counted argument uses, and every step of the counted proof built on it is retained verbatim in the delivered body.\par\medskip
\ifdefined\hypertarget\hypertarget{sourcepointer-2}{}\fi
\noindent\textit{Source reference.} Original reference 0FH5: sources/stacks-weil/pinned-original-source.tex, source lines 2233--2233. Prerequisite pointer invoked at proof line(s) 3706; it names a standard fact of the same chapter that the counted argument uses, and every step of the counted proof built on it is retained verbatim in the delivered body.\par\medskip
\ifdefined\hypertarget\hypertarget{sourcepointer-3}{}\fi
\noindent\textit{Source reference.} Original reference 0FHU: sources/stacks-weil/pinned-original-source.tex, source lines 3656--3656. Prerequisite pointer invoked at proof line(s) 3730; it names a standard fact of the same chapter that the counted argument uses, and every step of the counted proof built on it is retained verbatim in the delivered body.\par\medskip
\begin{lemma}
\label{lemma-classes-points}
Assume given (D0), (D1), (D2), and (D3) satisfying (A), (B), and (C).
Let $k'/k$ be a finite separable extension. Let $X$ be a geometrically
irreducible smooth projective scheme over $k'$ of dimension $d$.
Then $\gamma : \CH_0(X) \to H^{2d}(X)(d)$ factors through
$\deg : \CH_0(X) \to \mathbf{Z}$.
\end{lemma}

\begin{proof}
By Lemma \ref{lemma-generated-by-separable} it suffices to show: given
closed points $x, x' \in X$ whose residue fields are separable over $k$
we have $\deg(x') \gamma([x]) = \deg(x) \gamma([x'])$.

\medskip\noindent
We first reduce to the case of $k'$-rational points. Let $k''/k'$ be a
Galois extension such that $\kappa(x)$ and $\kappa(x')$ embed into $k''$
over $k$. Set $Y = X \times_{\Spec(k')} \Spec(k'')$ and denote $p : Y \to X$
the projection. By our choice of $k''/k'$ there exists a
$k''$-rational point $y$, resp.\ $y'$ on $Y$ mapping to $x$, resp.\ $x'$.
Then $p_*[y] = [k'' : \kappa(x)][x]$ and
$p_*[y'] = [k'' : \kappa(x')][x']$ in $\CH_0(X)$.
By compatibility with pushforwards given in axiom (C)(b)
it suffices to prove $\gamma([y]) = \gamma([y'])$ in $\CH^{2d}(Y)(d)$.
This reduces us to the discussion in the next paragraph.

\medskip\noindent
Assume $x$ and $x'$ are $k'$-rational points. By
Lemma \ref{lemma-divide-difference-points} there
exists a finite separable extension $k''/k'$ of fields
such that the pullback $[y] - [y']$
of the difference $[x] - [x']$ becomes divisible
by an integer $n > 1$ on $Y = X \times_{\Spec(k')} \Spec(k'')$.
(Note that $y, y' \in Y$ are $k''$-rational points.)
By Lemma \ref{lemma-relations-classes-points} we have
$\gamma([y]) = \gamma([y'])$ in $H^{2d}(Y)(d)$.
By compatibility with pushforward in axiom (C)(b)
we conclude the same for $x$ and $x'$.
\end{proof}

\end{document}
